% !TeX root = ../navier-stokes-manuscript.tex
% ============================================================
% 2.2 Finite-Time Response Escalation
% ============================================================

\subsection{Finite-Time Response Escalation}\label{sub:FiniteTimeResponseEscalation}

The Zeno schedule leaves open the only question that matters dynamically: can one solution perform all of those changes before the available time runs out?
A scale can shrink on paper without any fluid moving through the intervening states.
The Navier--Stokes field has no such freedom.
Every finer state must be reached from the state already present, and a finite endpoint shortens the time in which that response can occur.

Follow a particle of the same fluid.
Let \(X(a,t)\) solve
\[
  \frac{d}{dt}X(a,t)=u(X(a,t),t),
  \qquad
  X(a,0)=a,
\]
and write
\[
  D_t:=\partial_t+u\cdot\nabla
\]
for differentiation along that path.
Velocity, acceleration, and jerk are the first three rungs of one material response:
\[
  A_0(a,t):=u(X(a,t),t),
  \qquad
  A_1(a,t):=D_tu(X(a,t),t),
  \qquad
  A_2(a,t):=D_t^2u(X(a,t),t).
\]
More generally, set
\[
  A_r(a,t):=D_t^r u(X(a,t),t).
\]
The rungs are not independent because
\[
  \frac{d}{dt}A_r(a,t)=A_{r+1}(a,t).
\]
On each compact subinterval of \([0,T_*)\), classicality makes \(A_r(a,\cdot)\) absolutely continuous and gives
\[
  A_r(a,t)-A_r(a,s)
  =
  \int_s^t A_{r+1}(a,\tau)\,d\tau
\]
for \(0\le s<t<T_*\).
Consequently,
\[
  \lvert A_r(a,t)-A_r(a,s)\rvert
  \le
  \int_s^t \lvert A_{r+1}(a,\tau)\rvert\,d\tau.
\]
If \(A_r(a,t)\) is unbounded as \(t\uparrow T_*\), then
\[
  \int_0^{T_*}\lvert A_{r+1}(a,\tau)\rvert\,d\tau
  =
  \infty,
\]
and the finite length of \([0,T_*)\) forces
\[
  \sup_{0<t<T_*}\lvert A_{r+1}(a,t)\rvert
  =
  \infty.
\]
The escape of one material response forces escalation in the response immediately above it.

This statement is conditional, and that condition matters.
It does not say that every solution climbs the ladder, that any rung is monotone, or that every response is nonzero.
A steady solution can remain steady, and the zero solution can keep every rung equal to zero.
The ladder describes what finite-time escape would force inside the one fluid if escape occurred.

The same dependence appears before choosing a particle, but its functional setting must stay fixed throughout the implication.
Fix one Banach space \(B\), retain the Eulerian rows
\[
  V_n(t):=\partial_t^n u(\cdot,t),
\]
and suppose that \(V_n:[0,T)\to B\) is absolutely continuous on every compact subinterval, with Bochner derivative \(V_{n+1}\).
For \(0\le s<t<T<T_*\), the Banach-valued fundamental theorem of calculus gives
\[
  V_n(t)-V_n(s)
  =
  \int_s^t V_{n+1}(\tau)\,d\tau
\]
and
\[
  \lVert V_n(t)-V_n(s)\rVert_B
  \le
  \int_s^t \lVert V_{n+1}(\tau)\rVert_B\,d\tau.
\]
If
\[
  \sup_{0<t<T_*}\lVert V_n(t)\rVert_B
  =
  \infty,
\]
then the finiteness of \(\lVert V_n(0)\rVert_B\) forces
\[
  \int_0^{T_*}\lVert V_{n+1}(\tau)\rVert_B\,d\tau
  =
  \infty.
\]
Because \(T_*<\infty\), this in turn forces
\[
  \sup_{0<t<T_*}\lVert V_{n+1}(t)\rVert_B
  =
  \infty.
\]
The implication can be iterated upward only while the same fixed space \(B\), absolute continuity, and Bochner-derivative identity hold at every successive row.
That is the precise Eulerian form of finite-time response escalation.

Material and Eulerian responses still cannot be identified.
Already the material acceleration is
\[
  D_tu
  =
  \partial_tu+(u\cdot\nabla)u,
\]
and the material jerk expands to
\[
  D_t^2u
  =
  \partial_t^2u
  +2(u\cdot\nabla)\partial_tu
  +(\partial_tu\cdot\nabla)u
  +(u\cdot\nabla)\bigl((u\cdot\nabla)u\bigr).
\]
One step up in material response has already mixed a new time derivative with new spatial derivatives and nonlinear products of lower rows.
The particle picture exposes the escalation, while the equation decides how its pieces are coupled.

A single ladder indexed only by time can no longer hold what the fluid is doing.
The response has one index for temporal depth and another for spatial depth, and pressure enters as soon as the equation is used to relate them.
The next object must record that full mixed response without splitting it into separate derivative stories.
